\section{Metodología}

Se aplicó un procedimiento descriptivo y reproducible. El Documento Padre se utilizó como fuente académica principal; el archivo de 3.000 incidentes se conservó en \texttt{data/original}; y los eventos y encuestas necesarios para calcular las métricas se generaron mediante una simulación local con semilla 25022.

El trabajo se realizó en cuatro pasos: análisis y clasificación de incidentes; definición de tareas, atributos y métricas; generación y procesamiento automatizado de evidencia; e interpretación de KPI y formulación del plan de mejora. Python se utilizó para clasificación, simulación y cálculo; React y Flutter para representar las tareas de los usuarios; y los servicios locales para registrar los fallos controlados.

\begin{table}[H]
\centering
\caption{Fuentes empleadas en la evaluación}
\begin{tabularx}{\textwidth}{p{3cm}p{2.2cm}YY}
\toprule
Fuente & Cantidad & Procedencia & Uso \\
\midrule
Incidentes & 3.000 & Documento Padre & Clasificación ISO/IEC 25022 \\
Eventos & 6.000 & Simulación, semilla 25022 & Efectividad, eficiencia, riesgo y contexto \\
Encuestas & 2.000 & Simulación, semilla 25022 & Satisfacción y abandono \\
Catálogo & 10 métricas & Definición del equipo & Fórmulas, metas y semáforo \\
\bottomrule
\end{tabularx}
\end{table}

La integridad del dataset se verificó comparando el SHA-256 de la copia de trabajo con el archivo del Documento Padre. Ambos produjeron el mismo resultado:

\begin{center}
\footnotesize\ttfamily
814F3E5D8CF983A86BC125EC9C8461E3728C91DDCE1CFE8A5518ABD7A53CDC17
\end{center}

Los escenarios 1 a 3 siguen el contenido suministrado; los escenarios 4 a 12 se desarrollan a partir de su índice y se identifican como reconstruidos.
